\section{Antecedentes}

\subsection{Antecedentes internacionales}
\textcite{li2020} propusieron una red de segmentación semántica con estrategia de
aprendizaje auxiliar para delimitar celdas Braille, aunque su enfoque exige una
clasificación posterior. \textcite{ovodov2021} planteó el problema como detección
de objetos sobre la celda completa, lo que mejora la robustez ante rotaciones.
\textcite{sun2025} exploraron variantes de YOLOv11 para detectar Braille en
escenas naturales. En conjunto, estos trabajos respaldan la viabilidad técnica
del enfoque de detección de celda completa.

\subsection{Antecedentes nacionales}
Para el español, \textcite{fuentes2026} presentaron una aplicación móvil basada en
CNN para traducción bidireccional de Braille, pero no evalúan Braille manuscrito
escolar ni el compromiso precisión--eficiencia en hardware de gama baja.

\subsection{Síntesis y aporte diferencial}
La literatura comparada reporta mAP, F1, CER y latencia como métricas habituales.
Persisten tres vacíos: el Braille escolar manuscrito, la ejecución sin conexión en
dispositivos de gama media/baja y una evaluación integrada de precisión,
concordancia, eficiencia y efecto pedagógico. Este proyecto atiende los tres.
